\section{Background}
\subsection{Traditional Machine Learning Models}

This paper uses seven traditional machine learning models for the classification task: Random Forest, Support Vector Machine, Gradient Boosting, Decision Tree, $k$-Nearest Neighbors, Extreme Learning Machine (ELM), and Multilayer Perceptron (MLP). The rationale for selecting these models is explained in the Methodology section. Below is a brief introduction to each model.

\subsubsection{Random Forest}

Random Forest (RF) is an ensemble learning method first proposed by Ho in 1995 \cite{Ho1995} and later developed by Breiman and Cutler \cite{Breiman2001}. The algorithm grows multiple decision trees and combines their outputs to produce a final prediction. Each tree is trained using a random vector $\Theta$, generated independently from past vectors but with the same distribution \cite{Breiman2001}. The final prediction is obtained by aggregating ("voting") the predictions of all trees, typically via majority vote for classification \cite{Breiman2001}. RF is known for its robustness to overfitting, ability to handle high-dimensional data, and effectiveness in capturing complex interactions between features \cite{Breiman2001}. It remains one of the most effective machine learning models despite its early introduction \cite{Delgado2014}.

\subsubsection{Support Vector Machine}

Support Vector Machine (SVM) is a supervised learning model based on a simple idea: find an optimal hyperplane that separates the input data into different classes. The optimal hyperplane is defined as the one that maximizes the margin between classes, where the margin is the distance between the hyperplane and the nearest training points from each class (support vectors) \cite{Cortes1995}. SVM can be extended to handle non-linearly separable data by using kernel functions, particularly the Gaussian Radial Basis Function (RBF) kernel, which maps inputs into a higher-dimensional feature space where a linear separation is possible \cite{Schoelkopf2002}. SVM is widely used for classification tasks due to its effectiveness in high-dimensional spaces and its ability to model complex relationships between features and classes \cite{Hsu2003}.

\subsubsection{Gradient Boosting}

Gradient Boosting (GB) is an ensemble learning method that builds a series of weak learners, typically trees, that learn from the residuals of previous learners. Residuals are the differences between the predicted values and the actual values of the target variable \cite{Chen2024}. The training algorithm iterates through $M$ boosting rounds. In each round $m$, a new weak learner $T_{(m)}$ is trained to predict the residuals, thereby rectifying mistakes made by previous learners \cite{Chen2024}. The final prediction is obtained by summing the predictions of all weak learners, weighted by a learning rate $\nu$ that controls each learner's contribution \cite{Friedman2001}. The prediction process works as follows: for each tree $T_{(m)}$ and test instance $x_i$, the output $T_m(x_i)$ is scaled by $\nu$ and added to the cumulative prediction from previous trees, resulting in
\[
	F_M(x_i) = \sum_{m=1}^M \nu T_m(x_i). 
\]
GB is known for its high predictive accuracy and is widely used in various applications, including classification and regression tasks \cite{Friedman2001}.

\subsubsection{Decision Tree}

Decision Tree (DT) is a supervised learning model that recursively partitions the feature space into regions that are as homogeneous as possible with respect to the target label. At each internal node, the algorithm selects a feature (and a threshold for numerical features) that maximizes an impurity reduction criterion such as Gini impurity or information gain. Predictions are obtained by following the decision rules from the root to a leaf and outputting the majority class for classification tasks. Decision trees are interpretable and can capture non-linear decision boundaries, but they may overfit without appropriate regularization (e.g., maximum depth, minimum samples per split/leaf, or pruning).

\subsubsection{$k$-Nearest Neighbors}

$k$-Nearest Neighbors ($k$-NN) is a non-parametric, instance-based learning algorithm that classifies a data point based on the majority class among its $k$ nearest neighbors in feature space \cite{Cover1967}. Distance between data points is typically measured by Euclidean distance, although other measures such as Pearson or Spearman distance can also be used \cite{Jaskowiak2011}. The choice of the algorithm's sole hyperparameter, $k$, is crucial for performance. A larger value of $k$ can make the model more robust to noise, though it may blur interclass boundaries \cite{Everitt2011}. Due to its simplicity, $k$-NN's accuracy can be significantly affected by the choice of $k$ and the nature of the data, and it often requires feature scaling to ensure that all features contribute equally to the distance calculations \cite{Nigsch2006}. $k$-NN is noted for its strong consistency: as the dataset size increases, the probability of $k$-NN making an error converges to no worse than twice the Bayes error rate, which is the lowest possible error rate for any classifier \cite{Cover1967}.

\subsubsection{Extreme Learning Machine}

Extreme Learning Machine (ELM) is a type of feedforward neural network, first proposed by Huang et al. in 2006 \cite{Huang2006}. ELM is a learning algorithm for single-hidden layer feedforward neural networks (SLFNs) whose training is significantly faster than traditional counterparts. ELM randomly assigns input weights and hidden-layer biases, and then analytically determines the output weights by solving a linear system of equations \cite{Huang2006}. The training process consists of three steps: (1) randomly generate the input weights and hidden-layer biases; (2) compute the hidden-layer output matrix; (3) calculate the output weights using the Moore--Penrose generalized inverse of the hidden-layer output matrix \cite{Huang2006}. ELM has been shown to achieve good generalization performance with significantly reduced training time compared to traditional feedforward neural networks, making it suitable for large-scale classification tasks \cite{Huang2012}.

\subsubsection{Multilayer Perceptron}

Multilayer Perceptron (MLP) is a class of feedforward artificial neural network that consists of at least three layers: an input layer, one or more hidden layers, and an output layer \cite{Rosenblatt1958}. First proposed by Rosenblatt in 1958, MLP also serves as a basis for more complex neural network architectures, upon which ELM is built. Crucial in MLP architecture are the concepts of units and hidden layers. Units, also known as neurons, are the basic computational elements of MLP that receive input, process it, and produce output. Hidden layers are the layers of units that exist between the input and output layers, allowing MLP to learn complex representations of data \cite{Rosenblatt1958}. MLP functions by assigning weights to the connections between units and applying an activation function (commonly a sigmoid or ReLU function) to the output of each neuron. Learning in MLP is typically achieved through backpropagation and gradient descent algorithms, which adjust the weights to minimize the error between the predicted and actual outputs (modeled by a loss function) \cite{Haykin1999}.

\subsection{Large Language Models}

Large Language Models (LLMs) present a significant advancement in artificial intelligence, particularly in the field of natural language processing (NLP). The architecture of LLMs is based on the Transformer model, first introduced in a foundational paper by Vaswani et al. in 2017 \cite{Vaswani2017}. The Transformer architecture utilizes self-attention mechanisms to process input data, allowing it to capture long-range dependencies and contextual information effectively \cite{Vaswani2017}. LLMs are trained on enormous amounts of data and often possess billions of parameters, enabling them to generate human-like text and perform a wide range of language tasks with high accuracy \cite{Brown2020}. The training process of LLMs typically involves two stages: pre-training and fine-tuning. During pre-training, the model learns general language patterns from a large corpus of text, while fine-tuning involves further training the model on specific tasks or domains to enhance its performance in those areas \cite{Devlin2018}. LLMs have revolutionized the field in that they can generalize across a wide range of tasks without task-specific training, and they have been shown to outperform traditional machine learning models in various NLP tasks \cite{Brown2020}.

For our paper, we choose two models: \textbf{Qwen-2.5-3B} and \textbf{LLaMA-3-8B} (the suffix $X\text{B}$ of each model's name refers to the number of billions of parameters $X$). These models are publicly available on Hugging Face. Below is a brief introduction to each model.

\subsubsection{Qwen-2.5-3B}

Qwen-2.5-3B is a large language model developed by Alibaba, which is part of the Qwen series of models. The Qwen 2.5 series is designed to be more efficient than its predecessors, especially on coding and mathematical tasks. It has also been reported to enhance 
instruction-following, producing longer responses (beyond 8K tokens), interpreting structured inputs (e.g., tables), and generating structured outputs, particularly in JSON format. It is also more robust to variations in system prompts, which improves role-playing behavior and condition control in chatbot applications.

\subsubsection{LLaMA-3-8B}

LLaMA-3-8B is a large language model developed by Meta as part of the LLaMA series. It is designed to be efficient and capable across diverse language tasks. Architecturally, LLaMA 3 is an autoregressive language model built on an optimized Transformer backbone. Its tuned variants are aligned with human preferences for helpfulness and safety through supervised fine-tuning (SFT), followed by reinforcement learning from human feedback (RLHF).